\section{Extended Domain Notes for Future Main-Text Cuts}

This section is a deliberate landing zone for material that was either shortened in the first round of main-text compression or is likely to be moved out during the next round. It is intentionally organized by the same high-pressure regions identified in the second-round plan, but unlike a placeholder appendix, it already carries mechanism-centered prose that preserves the original argumentative intent of the main manuscript.

\subsection{Access-Scheduling Design Space Restored from the Main Paper}

The last page-budget pass also moved the access-scheduling design-space visual out of the main paper. That move was acceptable only because the core argument remained in prose: state access becomes a control problem when the runtime must decide what is visible, when a transition is legal, who owns the object now, and how disturbance changes the cost of the next action. The figure still matters as a compact reminder that observability, locality, ownership transfer, and recovery are not separate subtopics but one coupled design space.

\begin{figure*}[t]
\centering
\resizebox{0.94\textwidth}{!}{\begin{tikzpicture}[
  font=\footnotesize,
  layer/.style={
    draw=black!40,
    rounded corners=2pt,
    fill=black!1,
    minimum width=118mm,
    minimum height=12mm,
    inner sep=4pt
  },
  focus/.style={
    draw=black!45,
    rounded corners=2pt,
    align=left,
    text width=35mm,
    minimum height=9mm,
    inner sep=3pt
  },
  systems/.style={
    draw=black!30,
    rounded corners=2pt,
    align=left,
    text width=48mm,
    minimum height=9mm,
    inner sep=3pt,
    fill=white
  },
  lab/.style={font=\footnotesize\bfseries, align=right, text width=28mm},
  arr/.style={-{Stealth[length=1.9mm]}, line width=0.65pt, draw=black!60, shorten <=2pt, shorten >=2pt},
  feed/.style={-{Stealth[length=1.9mm]}, line width=0.65pt, densely dashed, draw=black!55, shorten <=2pt, shorten >=2pt}
]

\node[layer] (l1) at (4.4,2.3) {};
\node[layer] (l2) at (4.4,0.7) {};
\node[layer] (l3) at (4.4,-0.9) {};

\node[lab] at (-2.45,2.3) {Observability};
\node[focus, fill=blue!6] (f1) at (1.15,2.3) {queue exposure\\conflict density\\progress frontiers};
\node[systems] (s1) at (6.0,2.3) {Aurora, MillWheel, Naiad, Trill, Differential Dataflow};

\node[lab] at (-2.45,0.7) {Cost and scheduling};
\node[focus, fill=green!6] (f2) at (1.15,0.7) {topology, skew\\NUMA, ordering\\admission, migration};
\node[systems] (s2) at (6.0,0.7) {BriskStream, MorphStream, Megaphone, SOLA, ExeGPT};

\node[lab] at (-2.45,-0.9) {Recovery and elasticity};
\node[focus, fill=purple!6] (f3) at (1.15,-0.9) {checkpoint alignment\\ownership transfer\\replay boundary};
\node[systems] (s3) at (6.0,-0.9) {Spacker, Flink, SONIC, ORION, fast recovery};

\foreach \i in {1,2,3} {
  \draw[arr] (f\i.east) -- (s\i.west);
}

\draw[arr] (f1.south) -- (f2.north);
\draw[arr] (f2.south) -- (f3.north);
\draw[feed] (s3.south) -- ++(0,-0.38) -| (f1.west);

\end{tikzpicture}
}
\caption{Access-scheduling design space for locality, legality, transfer, and recovery.}
\label{fig:supp-access-designspace}
\end{figure*}

What should remain visible after the move is the comparative lesson behind the figure. Visibility mechanisms tell the runtime when a state transition is safe to expose; locality mechanisms tell it which placement keeps the next access affordable; ownership-transfer mechanisms tell it when relocation may occur without dual writers or replay ambiguity; recovery mechanisms tell it how those same promises survive disturbance. These are not four optional implementation concerns. They are the four places where access control becomes system-wide behavior rather than local bookkeeping.

The broader access-scheduling lineage is equally important for keeping the survey's boundary crisp. Classical consistency and concurrency systems such as TicToc, Percolator, and RACE show that timestamping, metadata layout, and remote coordination are already control surfaces over state, even before modern serving runtimes made the same pattern explicit~\cite{yu2016tictoc,peng2010percolator,zuo2021race}. Systems such as CheckFreq, ChecknRun, PipeFill, and the newer runtime-monitoring lines show that the access problem is not just placement, but also when the runtime may safely observe, checkpoint, and resume without invalidating later actions~\cite{mohan2021checkfreq,eisenman2022checknrun,arfeen2024pipefill,zhou2024neomem}. Tooling for dynamic reconfiguration and transfer, including Caerus, MORPH-style recovery, and queue-aware stream control, matters because it keeps the state-ownership question visible instead of treating movement as a mere implementation detail~\cite{zhang2021caerus,mao2023morphstream,fu2024autocheck,gao2025deck}.

The same lineage also includes systems that turn placement pressure into explicit runtime policy rather than hidden implementation cost. Resource-multiplexing and scheduler-centric systems such as G10, resource multiplexing, PFS-style reconfiguration, and workload-sensitive monitoring show that queueing, hot-state exposure, and placement are inseparable once the runtime manages shared state across stages or tenants~\cite{jiang2024megascale,he2025resourcemultiplexing,domingo2021pfsck,jiang2025shadowviews}. Transactional dataflow and distributed-storage systems such as Flash-style state flow, Cloud-style buffering, and multi-stage transfer control reinforce the same point: the access problem is always also a visibility problem and a handoff problem~\cite{fu2022sfs,fu2024autocheck,gao2025deck,li2020pegasus}. That is why the survey keeps access and scheduling as a first-class axis rather than folding it into generic runtime management.

This access line also needs a few more boundary references to keep the historical path intact. Transaction-processing systems such as TicToc, Percolator, and TiDB-style concurrency control show that control surfaces over mutable state have long existed in distributed data systems, even if they were not described in the same vocabulary as modern state management~\cite{yu2016tictoc,peng2010percolator}. Stream-processing systems such as Corun, Elastic stream systems, and BriskStream-style monitoring show that scheduling, state exposure, and reconfiguration were already intertwined before serving and retrieval brought those issues to the foreground~\cite{zhang2016corun,zhang2016elastic,zhang2019briskstream}. More recent monitoring and admission lines such as CheckFreq, ChecknRun, and PipeFill show that the runtime often needs an explicit contract for when observations can be trusted after disturbance~\cite{mohan2021checkfreq,eisenman2022checknrun,arfeen2024pipefill}.

The tail of the access line also includes a few systems that are easy to overlook because they sit between streaming, scheduling, and memory management. More recent scheduling and observability papers such as MOC, Optimum-style optimization, FS-MoE runtime control, and tree-aware transfer paths show that access control is still the main lever whenever a runtime has to reason about hidden contention and state movement together~\cite{cai2025moc,feng2025optimus,fsmoe25,wang2025wlbllm}. Likewise, Tigon, VM-control, and the updated memory-aware runtime lines keep the same lesson visible: once state crosses a device or tenancy boundary, legality and placement become one problem~\cite{huang2025tigon,zheng2024vmcu,tabatabai2024fbmm}. Even storage-to-compute bridges such as Snowflake-style disaggregation and mixed-access control belong in this lineage because they expose how latency, placement, and handoff interact~\cite{vuppalapati2020snowflake,zhou2024neomem}.

Two compact boundary references keep the access story complete. DRTMH-style transaction scheduling and the state-machine tradition remind us that legality over mutable state was a control problem long before modern runtimes named it that way~\cite{ren2015drtmh,schneider1990state}. EndMyth-style stream repair work shows the same issue from the failure side: access, recovery, and trust in observations are inseparable once the system can replay or resume~\cite{zamanian2017endmyth}.

\subsection{Serving-Lifecycle Notes}

The current main paper already compresses the serving literature into a shorter lifecycle narrative, but that compression necessarily removed several distinctions that still matter for auditability. The first distinction is between \emph{predictive}, \emph{structural-reuse}, and \emph{temporal-lifecycle} controllers. Predictive-serving systems such as WindServe, GLLM, Past-Future, HETIS, Seesaw, and ThunderServe do not all solve the same problem; they forecast different future debts and then bind those forecasts to different control surfaces~\cite{windserve2025,guo2025gllm,gong2025pastfuture,mo2025hetis,su2025seesaw,jiang2025thunderserve}. Some forecast phase-imbalance debt, some forecast future KV-memory peaks, and others forecast heterogeneous placement debt. The comparative lesson is not merely that prediction helps, but that a useful prediction must be tied to an explicit downstream debt model over queue growth, eviction risk, or transfer overhead.

The second distinction is between \emph{structural reuse} and \emph{footprint reduction}. Structural-reuse systems such as PromptCache, Marconi, ICCache, CacheBlend, CacheSlide, PIE, and speculative-prefix families make reuse legality visible by exposing prompt segmentation, compatibility, branch validity, or chunk-indexed reuse scope~\cite{gim2024promptcache,pan2025marconi,yu2025iccache,yao2025cacheblend,cacheslide2025,gim2025pie,miao2024specinfer}. Compression or paging systems, by contrast, mainly reduce the size or movement cost of state that is already known to be relevant. Keeping these families distinct is important because they answer different control questions: the former asks \emph{what may be reused safely}, while the latter asks \emph{how an already legal state object should be represented, placed, or transferred}.

The third distinction concerns \emph{where in time} the runtime intervenes. Pensieve, CachedAttention, HCache, Mooncake, and Symphony operate at the session-restoration boundary by preserving or restoring reuse value across turns~\cite{yu2025pensieve,gao2025cachedattention,gao2025hcache,qin2025mooncake,agarwal2026symphony}. FastServe, token-fair serving, LLMnix, MELL, and Libra operate at the active-contention boundary, where the central issue is how preemption, fairness, and migration distribute latency debt while requests compete live for the same state pool~\cite{wu2026fastserve,sheng2024fairnessserving,sun2024llumnix,liu2025mell,ruan2026libra}. HydraServe, ServerlessLLM, LayerKV, and Prefill-Only operate earlier at the bootstrap-residency boundary, showing that cold-start loading, tiered initialization, and early-layer retention already determine which state can ever become reusable later~\cite{lou2026hydraserve,fu2024serverlessllm,xiong2024layerkv,du2025prefillonly}. This time-structured comparison was compressed heavily in the main text, but it should remain preserved here because it explains why ``better caching'' is too weak a label for modern serving control.

Additional serving papers should be read as boundary refinements rather than as a separate taxonomy. SplitZip, CacheFlow, and ThunderServe make transfer cost and cloud heterogeneity explicit; ChandGPT-style reuse systems, Jenga, Prism, and Aegaeon make multi-model or multi-tenant residency explicit; Spinfer, FlexGen, and related offload designs make memory pressure and phase placement explicit; HeterInfer, MoE-Lightning, and resource multiplexing make the hardware/topology asymmetry explicit~\cite{guo2026splitzip,nian2026cacheflow,jiang2025thunderserve,gao2025weaver,zhang2025jenga,yu2025prism,fan2025spinfer,sheng2023flexgen,chen2025heteroinfer,cao2025moelightning,he2025resourcemultiplexing}. These works are useful here because they keep the lifecycle argument from shrinking into a single KV-cache story.

The serving cluster also needs a few bridge references that keep old and new state-governance lines connected. PagedAttention, Orca, and AlpaServe explain why batching and scheduling were already state-sensitive before the latest disaggregation wave; SlimPipe, What-If Stragglers, and Pumped trace-based systems explain why scheduling debt becomes visible only after phase or hardware asymmetry is modeled; output-length prediction and cache-aware routing systems show why admission and future-memory forecasting belong in the same discussion as restoration and transfer~\cite{kwon2023pagedattention,yu2022orca,li2023alpaserve,liu2025slimpipe,lin2025whatifstragglers,gong2025pastfuture,jiang2025shadowviews}. That bridge keeps the survey's lineage coherent without bloating the main paper.

Three more serving families deserve to stay visible as boundary references. Application-centric serving such as Parrot shows that workflow-level state and prompt DAGs can dominate placement decisions even before the cache story starts~\cite{lin2024parrot}. Tree-shaped speculative and group-wise systems such as FastTree and InstAttention show that reuse and scheduling depend on structure in the request stream, not just on raw throughput targets~\cite{pan2025fasttree,pan2025instattention}. Decentralized and overlay-based systems such as PlanetServe and GLLM-adjacent coordination lines show that cache reuse, routing, and trust become intertwined once the scheduler is no longer centralized~\cite{fang2026planetserve}. The same cluster also includes systems such as Infinity-style eviction, Duplex-style cross-turn reuse, and world-load balancing variants, which are best understood as refinements of lifecycle control rather than new top-level subfields~\cite{infinigen2024,yun2024duplex,wang2025wlbllm}.

Two final serving references are useful as compact examples of the same lifecycle logic. Prefill-optimized training-serving crossover systems such as MEPipe show that the serving boundary is often governed by how pipeline slack is distributed, not just by the cache itself~\cite{sun2025mepipe}. Copy-aware or clone-style state movement lines show the same thing from a different angle: once request state is shareable, the runtime has to decide whether cloning, relocating, or reusing is cheaper than recomputing~\cite{tian2025clone}.

The remaining serving papers that have not yet been mentioned directly still fit the same lifecycle view. Multi-model and adapter-aware systems such as Aegaeon, DLoRA, Chameleon, and Toppings expose tenant-locality and residency pressure; cache and reuse systems such as PromptCache, CacheWild, and Droidspeak expose the boundary between prompt overlap and safe reuse; inference-control systems such as Alise, Atom, and Prism-like lineage expose the boundary between active scheduling and future debt; and cloud-disaggregation systems such as Mooncake and PRFaaS expose the difference between local phase control and cross-cluster ownership~\cite{xiang2025aegaeon,wu2024dlora,iliakopoulou2024chameleon,li2025toppings,gim2024promptcache,wang2025kvcachewild,liu2026droidspeak,zhao2024alise,zhao2024atom,yu2025prism,qin2026prfaas}. Those distinctions are exactly what the compressed main paper needs the supplement to remember.

\begin{keybox}
The serving cluster should be read as one governed short-lived lifecycle: \emph{admit -> place -> share -> reclaim -> transfer -> restore}. First-round main-text compression preserved that thesis but necessarily removed much of the family-level detail explaining how different systems attach to different stages of that loop.
\end{keybox}

\begin{figure*}[t]
\centering
\resizebox{0.94\textwidth}{!}{\begin{tikzpicture}[
  font=\footnotesize,
  stage/.style={
    draw=black!45,
    rounded corners=2pt,
    align=center,
    text width=17mm,
    minimum height=10mm,
    inner sep=3pt
  },
  ctrl/.style={
    draw=black!30,
    rounded corners=2pt,
    align=center,
    text width=18mm,
    minimum height=8mm,
    inner sep=2pt,
    fill=black!2,
    font=\scriptsize
  },
  arr/.style={-{Stealth[length=1.7mm]}, line width=0.65pt, draw=black!65, shorten <=1pt, shorten >=1pt},
  feed/.style={-{Stealth[length=1.7mm]}, line width=0.65pt, densely dashed, draw=black!55, shorten <=1pt, shorten >=1pt},
  note/.style={font=\scriptsize, text=black!62, align=center}
]

\node[stage, fill=blue!6]   (s1) at (0,0) {\textbf{Admit}\\request};
\node[stage, fill=green!6]  (s2) at (2.25,0) {\textbf{Place}\\KV pages};
\node[stage, fill=yellow!10] (s3) at (4.5,0) {\textbf{Mutate}\\token state};
\node[stage, fill=purple!6] (s4) at (6.9,0) {\textbf{Expose /}\\\textbf{Transfer}};
\node[stage, fill=orange!8] (s5) at (9.3,0) {\textbf{Compact}\\fragments};
\node[stage, fill=red!5]    (s6) at (11.7,0) {\textbf{Evict /}\\\textbf{Reclaim}};

\foreach \a/\b in {s1/s2,s2/s3,s3/s4,s4/s5,s5/s6} {
  \draw[arr] (\a.east) -- (\b.west);
}

\node[ctrl] (c1) at (0,-1.15) {SLO gate\\admission};
\node[ctrl] (c2) at (2.25,-1.15) {residency\\plan};
\node[ctrl] (c3) at (4.5,-1.15) {growth and\\sharing edits};
\node[ctrl] (c4) at (6.9,-1.15) {reuse legality\\transfer};
\node[ctrl] (c5) at (9.3,-1.15) {fragmentation\\debt};
\node[ctrl] (c6) at (11.7,-1.15) {restore priority\\release};

\foreach \i in {1,2,3,4,5,6} {
  \draw[arr] (s\i.south) -- (c\i.north);
}

\draw[feed] (c6.south) -- ++(0,-0.35) -| (c1.south);
\node[note, text width=118mm] at (5.85,-2.25) {
  Memory pressure is not a terminal event: reclaiming KV state changes later restoration cost,
  prefix reuse, and tenant-visible SLO stability.
};

\node[note, text width=118mm] at (5.85,-2.85) {
  Representative systems: vLLM, Sarathi-Serve, SGLang, DistServe, Splitwise,
  Jenga, CacheFlow, Pensieve, Mooncake, FastServe.
};

\end{tikzpicture}
}
\caption{KV-cache lifecycle.}
\label{fig:supp-kvcache-lifecycle}
\end{figure*}

The lifecycle figure makes the control seams concrete. Admission decides whether a request is allowed to create new short-lived state at all. Placement decides where that state lives and therefore which queues or buses it will burden next. Sharing decides whether prefix, adapter, or session-local state may be reused safely. Reclamation decides what can be shrunk or evicted without hiding future restoration debt. Transfer and restoration decide whether the lost reuse value can be recovered cheaply enough to remain compatible with the original SLO.

If later rounds cut the main paper further, the following evidence should continue to accumulate in this subsection rather than disappear:
\begin{itemize}[leftmargin=*]
\item predictive-serving families grouped by the debt they forecast;
\item legality conditions for prefix, branch, and multi-call reuse;
\item temporal-lifecycle policies over restore, migrate, warm-start, and fairness debt;
\item clearer distinctions between allocator realism, transfer scheduling, and ownership semantics.
\end{itemize}

\subsection{Retrieval and Retention Governance Notes}

The first-round compression also shortened material whose original purpose was to keep retrieval and retention from reading like unrelated domains. The unifying point is that both expose a \emph{longer-lived memory layer} whose value decays gradually and whose maintenance debt is easy to hide if the runtime only reports steady-state quality. In retrieval systems, the state object is not just embeddings or an ANN graph in the abstract, but mutable shards, publication metadata, freshness signals, deletion debt, and planner-visible exposure state~\cite{singh2021freshdiskann,wang2021milvus,guo2022manu,asai2024selfrag,zhang2026flowrag,wang2026candor}. In retention systems, the state object is not just ``memory'' generically, but a bounded store whose protected parameters, exemplars, replay entries, compressed features, and effective budget must be co-governed over time~\cite{kirkpatrick2017ewc,rebuffi2017icarl,aljundi2019mir,prabhu2020gdumb,hayes2020remind,zhou2025ferret,li2026streamfp}.

Retrieval systems expose at least four trigger layers: request-time invocation, index-time maintenance, publication-time exposure, and placement-time residency. Self-RAG changes whether external memory should be consulted at all; FlowRAG changes when retriever state should be refreshed before reuse quality decays too far; CANDOR-Bench and FreshDiskANN expose when update churn, deletion debt, or local repair have already pushed the index away from its intended operating region; Milvus and Manu expose when partially refreshed state may be published safely under concurrent updates; Rummy and composable-memory search expose when movement or placement is justified by downstream quality gain instead of raw lookup throughput~\cite{asai2024selfrag,zhang2026flowrag,wang2026candor,singh2021freshdiskann,wang2021milvus,guo2022manu,zhang2024rummy,quinn2025accelerating}. The comparison shows that retrieval quality is a lifecycle property, not merely an encoder property.

Retention systems expose an analogous trigger hierarchy over one bounded store. EWC protects already-committed parameter state; iCaRL, GDumb, and StreamFP govern what enters the scarce retained set; GEM, A-GEM, MIR, and DER govern when stored state should re-enter the update path; REMIND governs whether retained signal is compressed and later reconstructed; Ferret governs what happens when the budget itself contracts under live update pressure~\cite{kirkpatrick2017ewc,rebuffi2017icarl,prabhu2020gdumb,li2026streamfp,lopez2017gem,chaudhry2019agem,aljundi2019mir,buzzega2020der,hayes2020remind,zhou2025ferret}. Read this way, the continual-learning literature matters to the survey not because it supplies another application area, but because it provides one of the clearest examples of how protect, admit, replay, compress, and budget-shrink decisions compose over time.

What should not be lost in later cuts is the comparative lesson across these two clusters. Retrieval runtimes lack a stable rule for when stale-but-local memory remains queryable; retention runtimes lack a stable rule for when compressed or demoted memory remains good enough to preserve. In both cases, the missing abstraction is a policy-visible exposure contract that turns hidden maintenance debt into an auditable control surface.

Several remaining papers fit this same long-lived-memory picture once the main paper's shorter synthesis is restored. Dynamic ANN and index-maintenance lines such as HNSW, DiskANN, Milvus, and graph- or shard-centric systems make clear that freshness, shard publication, and rebuild debt are all part of one operational budget~\cite{malkov2020hnsw,singh2021freshdiskann,wang2021milvus,guo2022manu}. Retrieval-augmented and memory-augmented systems such as DPR, ColBERT, REALM, RAG, ATLAS, HippoRAG, and Raptor show the same phenomenon at the model boundary: the runtime must manage what gets exposed, refreshed, abstracted, and reused rather than treating memory as a static auxiliary index~\cite{karpukhin2020dpr,khattab2020colbert,guu2020realm,lewis2020rag,izacard2022atlas,wang2024hipporag,sarthi2024raptor}. Those references are not extra background; they preserve the lifecycle logic behind the retrieval chapter.

Retention has a similar tail of boundary papers that the main paper can mention only briefly. Experience replay, memory budget, and continual-update lines such as MIR, GDumb, REMIND, Ferret, and StreamFP establish that protect/admit/replay/demote decisions are all budgeted governance acts over one bounded store~\cite{aljundi2019mir,prabhu2020gdumb,hayes2020remind,zhou2025ferret,li2026streamfp}. They are useful here because they let the survey connect long-horizon retention with short-lived serving reuse without pretending the two are the same problem.

\begin{figure*}[t]
\centering
\resizebox{0.94\textwidth}{!}{\begin{tikzpicture}[
  font=\footnotesize,
  head/.style={font=\scriptsize\bfseries, align=center, text=black!70},
  fam/.style={
    draw=black!45,
    rounded corners=2pt,
    align=center,
    text centered,
    text width=25mm,
    minimum height=10mm,
    inner sep=3pt,
    fill=blue!6
  },
  method/.style={
    draw=black!35,
    rounded corners=2pt,
    align=center,
    text centered,
    text width=42mm,
    minimum height=10mm,
    inner sep=3pt,
    fill=black!2
  },
  sys/.style={
    draw=black!35,
    rounded corners=2pt,
    align=center,
    text centered,
    text width=35mm,
    minimum height=10mm,
    inner sep=3pt,
    fill=white
  },
  arr/.style={-{Stealth[length=1.7mm]}, line width=0.64pt, draw=black!60, shorten <=2pt, shorten >=2pt},
  gap/.style={draw=red!35, fill=red!4, rounded corners=2pt, align=left, text width=104mm, inner sep=4pt, font=\scriptsize}
]

\node[head] at (0,0.95) {Strategy family};
\node[head] at (5.35,0.95) {Mechanism};
\node[head] at (10.7,0.95) {Representative systems};
\draw[black!45, line width=0.6pt] (-1.85,0.62) -- (13.0,0.62);

\node[fam, fill=blue!6] (f1) at (0,0) {\textbf{Protection}\\constrain drift};
\node[method] (m1) at (5.35,0) {importance weights\\gradient regularization};
\node[sys] (s1) at (10.7,0) {EWC, Online EWC,\\SI, MAS};

\node[fam, fill=green!6] (f2) at (0,-1.45) {\textbf{Replay}\\absorb interference};
\node[method] (m2) at (5.35,-1.45) {gradient projection\\reservoir sampling\\memory replay};
\node[sys] (s2) at (10.7,-1.45) {GEM, A-GEM, MIR,\\ER, DER++};

\node[fam, fill=yellow!10] (f3) at (0,-2.9) {\textbf{Admission}\\bound the store};
\node[method] (m3) at (5.35,-2.9) {exemplar anchoring\\utility scoring\\coreset selection};
\node[sys] (s3) at (10.7,-2.9) {iCaRL, GDumb,\\StreamFP, BiC};

\node[fam, fill=purple!6] (f4) at (0,-4.35) {\textbf{Compression}\\contract budget};
\node[method] (m4) at (5.35,-4.35) {feature distillation\\elastic budgets\\compressed replay};
\node[sys] (s4) at (10.7,-4.35) {REMIND, Ferret,\\PackNet, HAT};

\foreach \i in {1,2,3,4} {
  \draw[arr] (f\i.east) -- (m\i.west);
  \draw[arr] (m\i.east) -- (s\i.west);
}
\foreach \a/\b in {f1/f2,f2/f3,f3/f4} {
  \draw[arr] (\a.south) -- (\b.north);
}

\node[gap] at (5.35,-5.55) {
  \textbf{Open contract.}
  Admission, replay, compression, and demotion need one retention-store API
  so budget pressure can be compared with future reuse value.
};

\end{tikzpicture}
}
\caption{Retention-governance taxonomy.}
\label{fig:supp-retention-taxonomy}
\end{figure*}

The taxonomy figure matters because it prevents a recurring misreading: retention is not one replay heuristic with a few engineering variants. It is a store-level governance problem over multiple actions with different temporal meanings. Protect decisions stabilize already-committed state; admission decisions determine future coverage; replay decisions determine when retained state becomes active again; compression and demotion decisions decide how much of that future value survives under tighter budgets; budget-shrink rules decide which guarantees fail first when the store contract is stressed. That fuller action map was too expensive to keep in the main paper, but it should remain explicit here.

\subsection{Long-Horizon Evolution Anchors}

The longest-horizon evidence in the repository also belongs here because it keeps the survey from collapsing ``evolution'' into a narrow continual-learning story. SuperNeurons and related memory-reuse systems show that activation liveness, recomputation, and offload policy already form a managed state loop in deep-learning training, not just in inference-side serving~\cite{wang2018superneurons}. Bubble-filling and recomputation-aware training systems such as Obscura show the same thing from another angle: when state is expensive to keep live, the runtime must decide which debt to shift into pipeline slack and which debt to pay immediately~\cite{huang2025obscura}. That is the same governance pattern that later appears in retention stores, but with a different temporal horizon and a different type of reuse value.

\subsection{Blueprint and Fault-Model Notes}

The first piece that should remain visible here is the distinction between \emph{typed state objects}, \emph{observation contracts}, \emph{decision precedence}, \emph{actuation scope}, and \emph{safety-envelope invariants}. The blueprint is not a proposal for one monolithic runtime; its role is to make recurring hidden assumptions explicit enough that neighboring mechanisms can compose without silently violating one another's boundaries.

One implication is that recovery should not be treated as a separate semantic regime. If steady-state control reasons over typed KV pages, refreshable vector shards, replay buffers, or demotable retention objects, then recovery must restore and validate the same objects under the same exposure and action boundaries. Otherwise a system effectively runs one state model in steady state and another during failure. That inconsistency matters because it explains why fault-model notes are part of the mechanism, not optional decoration.

The fault-model distinction itself also deserves fuller preservation here than in the compressed main text. A safety claim that holds under crash-stop may fail under crash-recovery with replay, may fail differently under network partition, and may fail silently under degradation modes such as stale telemetry, compaction lag, or sudden budget contraction~\cite{chandy1985snapshot,corbett2012spanner,ongaro2014raft,mao2023morphstream,zhao2024recovery,wang2026candor,zhang2026flowrag}. These are not small implementation variants. They determine whether a controller may expose partially refreshed retrieval state, whether ownership may transfer across disaggregated serving pools, and whether a retention tier may demote memory before replaying it. The blueprint's value lies in precisely these boundary conditions.

Two invariant families deserve to be carried here explicitly if later main-paper cuts continue. The first is \emph{freshness exposure invariants}: a retrieval shard or index segment should specify when it is queryable for internal warmup, restricted traffic, or general external service after rebuild, deletion repair, or compaction. The second is \emph{retention demotion invariants}: a retention object should specify whether it is protected, compressible, summarizable, demotable, or droppable when effective budget shrinks. These invariants preserve the blueprint's operational meaning.

\subsection{Expanded Foundations Material}

The survey's inclusion rule is not venue prestige or application relevance. A paper is central when it turns state into an explicit runtime-managed object and exposes a reusable control surface over that state. That is why transactional streaming, migration, KV-cache lifecycle control, retrieval maintenance, and retention governance sit in the core, while many application-facing LLM, agent, or prediction papers stay outside it even when they are technically strong.

That definition also has a formal lineage. The state-machine tradition is useful here not because the survey is about formal methods, but because it reminds us that a system only becomes governable when the runtime can name, observe, and transform its state explicitly~\cite{schneider1990state}. That conceptual boundary is the reason the supplement keeps the access, serving, retrieval, and retention lines together instead of splitting them into unrelated application silos.

If later readers need a fuller rationale for why certain papers are core versus contextual, this supplement can further expand:
\begin{itemize}[leftmargin=*]
\item related-survey positioning,
\item broader literature landscape,
\item analytical boundary notes derived from repository triage decisions,
\item explicit mappings from triage labels to the five-field frame used in the main manuscript.
\end{itemize}

Two visuals removed in the final page-budget pass also belong conceptually with this foundations rationale. The first is the poster-style overview, which summarizes the survey spine and helps later readers see how the three axes and cross-domain synthesis fit together. The second is the running-example triptych, which explains why the paper reuses one access case, one serving case, and one retrieval case across sections instead of multiplying disconnected examples.

\begin{figure}[tbp]
\centering
\resizebox{0.9\linewidth}{!}{\begin{tikzpicture}[
  font=\footnotesize,
  box/.style={
    draw=black!38,
    rounded corners=2pt,
    align=center,
    text centered,
    minimum height=10mm,
    inner sep=3pt,
    text width=28mm,
    fill=white
  },
  dim/.style={box, text width=27mm, fill=blue!6},
  outbox/.style={box, text width=27mm, fill=green!5},
  group/.style={draw=black!25, rounded corners=2pt, inner sep=7pt},
  head/.style={font=\scriptsize\bfseries, align=center, text=black!70},
  arr/.style={-{Stealth[length=1.8mm]}, line width=0.65pt, draw=black!58, shorten <=2pt, shorten >=2pt},
  link/.style={line width=0.65pt, draw=black!58},
  conn/.style={-{Stealth[length=1.8mm]}, line width=0.65pt, draw=black!58, shorten >=2pt},
  feed/.style={line width=0.65pt, draw=black!50, dash pattern=on 4pt off 4pt, line cap=rect},
  feedarr/.style={-{Stealth[length=1.8mm]}, line width=0.65pt, draw=black!50, dash pattern=on 4pt off 4pt, line cap=rect, shorten <=2pt, shorten >=2pt}
]

\node[head] at (0,3.35) {System domains};
\node[box, fill=blue!4] (d1) at (0,2.20) {Streaming and\\transactional};
\node[box, fill=green!4] (d2) at (0,0.85) {Hardware-conscious\\execution};
\node[box, fill=yellow!8] (d3) at (0,-0.50) {LLM serving and\\inference};
\node[box, fill=purple!4] (d4) at (0,-1.85) {Retrieval and\\agentic memory};
\node[group, fit=(d1)(d2)(d3)(d4)] {};

\node[head] at (5.0,3.35) {Core dimensions};
\node[dim] (c1) at (5.0,1.75) {Access and\\scheduling};
\node[dim] (c2) at (5.0,0.20) {Execution\\optimization};
\node[dim] (c3) at (5.0,-1.35) {Evolution and\\reuse};
\node[group, fit=(c1)(c2)(c3)] {};

\node[head] at (9.8,3.35) {Survey outputs};
\node[outbox] (o1) at (9.8,1.75) {Comparative\\synthesis};
\node[outbox] (o2) at (9.8,0.20) {Runtime\\blueprint};
\node[outbox] (o3) at (9.8,-1.35) {Research\\agenda};
\node[group, fit=(o1)(o2)(o3)] {};

\coordinate (elbowRef) at ($(d1.east)!0.5!(c1.west)$);
\coordinate (j1) at ($(elbowRef |- c1.west)$);
\coordinate (j2) at ($(elbowRef |- c2.west)$);
\coordinate (j3) at ($(elbowRef |- c3.west)$);
\coordinate (feedX) at ($(j1)!0.68!(c1.west)$);
\coordinate (m1a) at ($(j1 |- d1.east)$);
\coordinate (m1b) at ($(j1 |- d2.east)$);
\coordinate (m2) at ($(j2 |- d3.east)$);
\coordinate (m3) at ($(j3 |- d4.east)$);
\coordinate (feedBottom) at ([yshift=-0.45cm]feedX |- j1);
\coordinate (feedTurn) at (feedX |- j1);

\draw[feed] (c3.south) -- ++(0,-0.45) -| (feedBottom);
\draw[feed] (feedBottom) -- (feedTurn);
\draw[feedarr] (feedTurn) -- (c1.west);

\draw[conn] (d1.east) -- (m1a) -- (j1) -- (c1.west);
\draw[link] (d2.east) -- (m1b) -- (j1);

\draw[conn] (d3.east) -- (m2) -- (j2) -- (c2.west);

\draw[conn] (d4.east) -- (m3) -- (j3) -- (c3.west);

\draw[arr] (c1.south) -- (c2.north);
\draw[arr] (c2.south) -- (c3.north);

\draw[arr] (c1.east) -- (o1.west);
\draw[arr] (c2.east) -- (o2.west);
\draw[arr] (c3.east) -- (o3.west);

\node[font=\scriptsize, text=black!58, align=center] at (5.0,-2.55)
  {access, execution, and evolution form one feedback loop};

\end{tikzpicture}
}
\caption{Poster-style survey overview restored in the supplement. It is no longer needed in the compressed main paper, but it remains useful as a compact map of the survey's scope and argumentative flow.}
\label{fig:supp-poster-overview}
\end{figure}

\begin{figure}[tbp]
\centering
\resizebox{0.96\linewidth}{!}{\begin{tikzpicture}[
  font=\footnotesize,
  cell/.style={
    draw=black!35,
    rounded corners=2pt,
    align=left,
    text width=27mm,
    minimum height=12mm,
    inner sep=3pt
  },
  rowlab/.style={font=\scriptsize\bfseries, align=right, text width=12mm, text=black!70},
  head/.style={font=\scriptsize\bfseries, align=center, text=black!72},
  arr/.style={-{Stealth[length=1.7mm]}, line width=0.6pt, draw=black!58, shorten <=2pt, shorten >=2pt}
]

\node[head] at (0,2.45) {State object};
\node[head] at (3.65,2.45) {Stress event};
\node[head] at (7.30,2.45) {Runtime decision};
\node[head] at (10.95,2.45) {Service effect};
\draw[black!45, line width=0.6pt] (-2.25,2.12) -- (12.55,2.12);

\node[rowlab] at (-2.15,1.05) {A};
\node[cell, fill=blue!4] (a1) at (0,1.05) {streaming fraud\\hot keys};
\node[cell, fill=blue!4] (a2) at (3.65,1.05) {campaign burst\\flips popularity};
\node[cell, fill=blue!4] (a3) at (7.30,1.05) {repartition hotspots\\adjust snapshots};
\node[cell, fill=blue!4] (a4) at (10.95,1.05) {tail latency recovers\\replay debt bounded};

\node[rowlab] at (-2.15,-0.25) {B};
\node[cell, fill=green!4] (b1) at (0,-0.25) {LLM assistant\\KV pages, adapters};
\node[cell, fill=green!4] (b2) at (3.65,-0.25) {long prompts create\\memory pressure};
\node[cell, fill=green!4] (b3) at (7.30,-0.25) {phase-aware admit\\selective transfer};
\node[cell, fill=green!4] (b4) at (10.95,-0.25) {TTFT stabilizes\\goodput survives};

\node[rowlab] at (-2.15,-1.55) {C};
\node[cell, fill=purple!4] (c1) at (0,-1.55) {RAG knowledge\\index, refresh queue};
\node[cell, fill=purple!4] (c2) at (3.65,-1.55) {corpus drift and\\delete backlog};
\node[cell, fill=purple!4] (c3) at (7.30,-1.55) {gate stale shards\\stage refresh};
\node[cell, fill=purple!4] (c4) at (10.95,-1.55) {freshness improves\\compaction controlled};

\foreach \r in {a,b,c} {
  \draw[arr] (\r1.east) -- (\r2.west);
  \draw[arr] (\r2.east) -- (\r3.west);
  \draw[arr] (\r3.east) -- (\r4.west);
}

\end{tikzpicture}
}
\caption{Running examples restored in the supplement. The main paper keeps the examples in prose; the supplement preserves the visual cue showing why the same three scenarios recur across access, execution, and evolution.}
\label{fig:supp-running-examples}
\end{figure}

\begin{keybox}
The supplement should evolve as a \emph{structured evidence archive}, not as a prose graveyard. Every future main-text cut should have an explicit landing place here before the cut is finalized.
\end{keybox}
